\section{必装插件}

作者推荐了四个 Claude Cowork 插件，按优先级排列。

\subsection{Productivity Plugin（优先安装）}

这是整个 Cowork 环境的基础，提供任务管理、日程安排和工作流自动化的 slash commands：

\begin{itemize}
    \item \texttt{/task} --- 在对话中直接创建和跟踪任务
    \item \texttt{/schedule} --- 直接在日历中安排时间块
    \item \texttt{/workflow} --- 运行已保存的多步骤自动化流程
\end{itemize}

\begin{knowledgebox}{为什么要优先安装}
Productivity Plugin 是其他插件的基础设施。后续的 Marketing、Data、Sales 等插件的工作流功能都依赖于它提供的 \texttt{/workflow} 命令。安装顺序影响整体体验。
\end{knowledgebox}

\subsection{Marketing Plugin}

面向内容创作者的核心插件，涵盖起草、多平台分发和内容日历规划：

\begin{itemize}
    \item \texttt{/draft-content} --- 起草 LinkedIn 帖子、推文、Newsletter 段落
    \item \texttt{/repurpose} --- 将一篇内容转化为 5 个平台的适配版本
    \item \texttt{/campaign} --- 规划包含主题和 hooks 的完整内容日历
\end{itemize}

作者特别指出，当 Marketing Plugin 与 brand voice 文件（见第~\ref{sec:brand-voice}~节）配合使用时，输出质量会显著提升。

\subsection{Data Plugin}

面向数据分析场景，将 Cowork 连接到实际数据源：

\begin{itemize}
    \item \texttt{/analyze} --- 上传 CSV 文件，无需写公式即可获取洞察
    \item \texttt{/build-dashboard} --- 创建交互式可视化看板
    \item \texttt{/sql} --- 用自然语言描述生成 SQL 查询
\end{itemize}

作者每周用该插件完成收入跟踪和客户报告。

\subsection{Sales Plugin（可选）}

适用于有外联需求的用户，将 30 分钟的会前准备压缩到 3 分钟：

\begin{itemize}
    \item \texttt{/research-account} --- 拉取公司信息、近期新闻、关键联系人
    \item \texttt{/prep-call} --- 在会议前生成简报文档
    \item \texttt{/draft-outreach} --- 基于潜在客户数据生成个性化邮件
\end{itemize}

\subsection{本章小结}

四个插件按优先级为：Productivity（基础设施）$\to$ Marketing（内容创作）$\to$ Data（数据分析）$\to$ Sales（销售外联）。其中 Productivity Plugin 是必装项，其余按需选择。

